% ═══════════════════════════════════════════════════════════════
% MUSTERLÖSUNG - Physik Klasse 6: Geschwindigkeit
% 25 BE (★) | 28 BE (★★) | 31 BE (★★★)
% ═══════════════════════════════════════════════════════════════

\documentclass[11pt, a4paper]{article}

\usepackage[utf8]{inputenc}
\usepackage[T1]{fontenc}
\usepackage[ngerman]{babel}
\usepackage{helvet}
\renewcommand{\familydefault}{\sfdefault}

\usepackage[margin=1.8cm, top=2cm, bottom=1.5cm]{geometry}
\usepackage{enumitem}
\usepackage{tabularx}
\usepackage{array}
\usepackage{colortbl}
\usepackage{tikz}
\usepackage{amssymb}
\usepackage{amsmath}
\usepackage{fancyhdr}
\usepackage{xcolor}
\usepackage{tcolorbox}
\tcbuselibrary{skins, breakable}
\usepackage{qrcode}
\usepackage[hidelinks]{hyperref}
\usepackage{siunitx}
\sisetup{
    locale = DE,
    per-mode = symbol,
    unit-font-command = \textsf
}
% Einheiten-Makros für einfache Verwendung
\DeclareSIUnit{\mps}{\meter\per\second}

% Lernpfad-URL
\newcommand{\lpurl}{https://devbox42.github.io/sciencesim/physik/kl06/02-bewegung/std-02-geschwindigkeit/LP-02-geschwindigkeit.html}

% Farben
\definecolor{physik}{HTML}{2c5aa0}
\definecolor{physik-hell}{HTML}{e8f0fa}
\definecolor{hellgrau}{HTML}{F5F5F5}
\definecolor{orange-warn}{HTML}{b35c00}
\definecolor{orange-hell}{HTML}{fff3e0}
\definecolor{loesung}{HTML}{c62828}

% Kopfzeile
\pagestyle{fancy}
\fancyhf{}
\fancyhead[L]{\small\textcolor{physik}{\textbf{Physik Klasse 6}} | \textcolor{loesung}{\textbf{MUSTERLÖSUNG}}: Geschwindigkeit}
\fancyhead[R]{\small\textcolor{loesung}{Nicht für Schüler!}}
\fancyfoot[R]{\small\thepage}
\renewcommand{\headrulewidth}{0pt}

% Boxen
\newtcolorbox{merkkasten}[1][]{
    colback=physik-hell,
    colframe=physik,
    boxrule=1.5pt,
    arc=0pt,
    left=8pt, right=8pt, top=8pt, bottom=6pt,
    title={\small\textbf{#1}},
    coltitle=white,
    fonttitle=\sffamily,
    attach boxed title to top left={yshift=-2mm, xshift=5mm},
    boxed title style={colback=physik, arc=0pt}
}

\newtcolorbox{aufgabenbox}[1][]{
    colback=white,
    colframe=physik,
    boxrule=1pt,
    arc=0pt,
    left=8pt, right=8pt, top=6pt, bottom=6pt,
    title={\small\textbf{#1}},
    coltitle=white,
    fonttitle=\sffamily,
    colbacktitle=physik,
    boxed title style={arc=0pt}
}

\newtcolorbox{protokollbox}{
    colback=orange-hell,
    colframe=orange-warn,
    boxrule=1.5pt,
    arc=0pt,
    left=8pt, right=8pt, top=8pt, bottom=6pt
}

% Befehle
\newcommand{\punkte}[1]{\hfill\textcolor{gray}{\small /#1}}
\newcommand{\lsg}[1]{\textcolor{loesung}{\textbf{#1}}}
\newcommand{\sternchen}[1]{\textcolor{physik}{\textbf{(#1)}}}

% ISO 80000: Formelzeichen kursiv/Serif, Einheiten aufrecht/Sans
\newcommand{\fz}[1]{\textit{\textrm{#1}}}           % Formelzeichen
\newcommand{\einheit}[1]{\textsf{\textup{#1}}}      % Einheit: aufrecht + Sans-Serif
\newcommand{\mps}{$\tfrac{\textsf{m}}{\textsf{s}}$} % m/s als Bruch

\setlength{\parindent}{0pt}
\setlength{\parskip}{4pt}

\begin{document}

% ═══════════════════════════════════════════════════════════════
% KOPFBEREICH
% ═══════════════════════════════════════════════════════════════

\noindent
\begin{minipage}[t]{0.82\textwidth}
\textcolor{physik}{\rule{\textwidth}{2pt}}\\[0.3em]
{\LARGE\bfseries\textcolor{physik}{Geschwindigkeit messen}}\\[0.2em]
\textcolor{loesung}{\large\textbf{--- MUSTERLÖSUNG ---}}\\[0.2em]
\textcolor{physik}{\rule{\textwidth}{0.5pt}}
\end{minipage}%
\hfill
\begin{minipage}[t]{0.15\textwidth}
\raggedleft
\href{\lpurl}{\qrcode[height=1.8cm]{\lpurl}}\\[-2pt]
{\tiny\textcolor{gray}{Lernpfad}}
\end{minipage}

\vspace{-0.3em}

\begin{flushright}
\small
\begin{tabular}{|c|c|c|c|}
\hline
\cellcolor{physik-hell} \textbf{$\star$ Basis} & \cellcolor{physik-hell} \textbf{$\star\star$} & \cellcolor{physik-hell} \textbf{$\star\star\star$} & \cellcolor{physik-hell} \textbf{Erreicht} \\
\hline
/25 BE & /28 BE & /31 BE & \hspace{1.5cm} \\
\hline
\end{tabular}
\end{flushright}

\vspace{0.5em}

% ═══════════════════════════════════════════════════════════════
% AUFGABE 1: Wiederholung Bewegungsarten (4 BE)
% ═══════════════════════════════════════════════════════════════

\begin{aufgabenbox}[Aufgabe 1: Bewegungsarten \sternchen{$\star$} \punkte{4}]
\textbf{Ordne zu:} Kreuze für jedes Beispiel die richtige Bahnform und Bewegungsart an.

\vspace{0.5em}

\renewcommand{\arraystretch}{1.6}
\begin{tabularx}{\textwidth}{|X|c|c|c|c|}
\hline
\rowcolor{physik-hell}
\textbf{Beispiel} & \textbf{geradlinig} & \textbf{krummlinig} & \textbf{gleichförmig} & \textbf{ungleichförmig} \\
\hline
a) Rolltreppe (konstant) & \lsg{$\boxtimes$} & $\square$ & \lsg{$\boxtimes$} & $\square$ \\
\hline
b) Karussell (dreht sich) & $\square$ & \lsg{$\boxtimes$} & \lsg{$\boxtimes$} & $\square$ \\
\hline
c) Auto bremst & \lsg{$\boxtimes$} & $\square$ & $\square$ & \lsg{$\boxtimes$} \\
\hline
d) Stein fällt & \lsg{$\boxtimes$} & $\square$ & $\square$ & \lsg{$\boxtimes$} \\
\hline
\end{tabularx}

\vspace{0.3em}
\textcolor{loesung}{\small\textit{Je Zeile 2 Kreuze = 1 BE (0,5 pro richtigem Kreuz)}}
\end{aufgabenbox}

\vspace{0.8em}

% ═══════════════════════════════════════════════════════════════
% AUFGABE 2: Experiment-Protokoll (11 BE)
% ═══════════════════════════════════════════════════════════════

\begin{protokollbox}
\textbf{\textcolor{orange-warn}{Aufgabe 2: Experiment – 100-Meter-Lauf \sternchen{$\star$} \punkte{11}}}

\vspace{0.5em}

\textbf{a) Messwerte:} Trage die Zeiten aus der Simulation ein und berechne die Geschwindigkeit.

\vspace{0.5em}

\renewcommand{\arraystretch}{2.2}
\begin{tabularx}{\textwidth}{|p{2.8cm}|c|>{\centering\arraybackslash}p{2.5cm}|>{\centering\arraybackslash}X|}
\hline
\rowcolor{physik-hell}
\textbf{Läufer} & \textbf{Strecke \fz{s}} & \textbf{Zeit \fz{t}} (messen!) & \textbf{Geschwindigkeit \fz{v}} (berechnen!) \\
\hline
Usain (Sprinter) & 100\,\einheit{m} & \lsg{10}\,\einheit{s} & \lsg{10}\,\mps \\
\hline
Lisa (Schülerin) & 100\,\einheit{m} & \lsg{12,5}\,\einheit{s} & \lsg{8}\,\mps \\
\hline
Tom (Jogger) & 100\,\einheit{m} & \lsg{20}\,\einheit{s} & \lsg{5}\,\mps \\
\hline
Oma Erna & 100\,\einheit{m} & \lsg{50}\,\einheit{s} & \lsg{2}\,\mps \\
\hline
\end{tabularx}

\vspace{0.3em}
\textcolor{loesung}{\small\textit{4 Zeiten (je 1 BE) + 4 Geschwindigkeiten (je 1 BE) = 8 BE}}

\vspace{0.8em}

\textbf{b) Merksatz:} Ergänze die Lücken.

\vspace{0.3em}

\begin{merkkasten}[Merksatz]
Die \textbf{Geschwindigkeit} $v$ gibt an, welchen \lsg{Weg} ein Körper

in einer bestimmten \lsg{Zeit} zurücklegt.

\vspace{0.8em}

\begin{center}
{\Large $v = \dfrac{\text{\lsg{$s$}}}{\text{\lsg{$t$}}}$}
\end{center}

\vspace{0.5em}

\textbf{Einheiten:} \fz{v} in \lsg{\mps}, \quad \fz{s} in \lsg{\einheit{m}}, \quad \fz{t} in \lsg{\einheit{s}}
\end{merkkasten}

\textcolor{loesung}{\small\textit{6 Lücken = 3 BE (je 0,5 BE)}}
\end{protokollbox}

\vspace{0.8em}

% ═══════════════════════════════════════════════════════════════
% AUFGABE 3: Berechnungen (6 BE)
% ═══════════════════════════════════════════════════════════════

\begin{aufgabenbox}[Aufgabe 3: Geschwindigkeit berechnen \sternchen{$\star$} \punkte{6}]
\textbf{Berechne} die Geschwindigkeit. Schreibe den vollständigen Rechenweg auf.

\vspace{0.5em}

\textbf{a)} Ein Radfahrer legt 200\,\einheit{m} in 40\,\einheit{s} zurück.

\vspace{0.3em}
\begin{tabular}{@{}ll}
Gegeben: & \fz{s} = \lsg{200\,\einheit{m}}, \quad \fz{t} = \lsg{40\,\einheit{s}} \\[0.5em]
Gesucht: & \lsg{\fz{v}} \\[0.5em]
Formel: & \fz{v} = \lsg{\fz{s} / \fz{t}} \\[0.5em]
Einsetzen: & \fz{v} = \lsg{200\,\einheit{m} / 40\,\einheit{s}} \\[0.5em]
Ergebnis: & \fz{v} = \lsg{5\,\mps}
\end{tabular}

\vspace{0.3em}
\textcolor{loesung}{\small\textit{Geg+Ges 1 BE, Formel 1 BE, Ergebnis 1 BE = 3 BE}}

\vspace{1em}

\textbf{b)} Eine Schnecke kriecht 6\,\einheit{m} in 600\,\einheit{s}.

\vspace{0.3em}
\begin{tabular}{@{}ll}
Gegeben: & \fz{s} = \lsg{6\,\einheit{m}}, \quad \fz{t} = \lsg{600\,\einheit{s}} \\[0.5em]
Gesucht: & \lsg{\fz{v}} \\[0.5em]
Formel: & \fz{v} = \lsg{\fz{s} / \fz{t}} \\[0.5em]
Einsetzen: & \fz{v} = \lsg{6\,\einheit{m} / 600\,\einheit{s}} \\[0.5em]
Ergebnis: & \fz{v} = \lsg{0,01\,\mps}
\end{tabular}

\vspace{0.3em}
\textcolor{loesung}{\small\textit{Geg+Ges 1 BE, Formel 1 BE, Ergebnis 1 BE = 3 BE}}
\end{aufgabenbox}

\vspace{0.8em}

% ═══════════════════════════════════════════════════════════════
% AUFGABE 4: Textaufgabe (4 BE)
% ═══════════════════════════════════════════════════════════════

\begin{aufgabenbox}[Aufgabe 4: Schulweg \sternchen{$\star$} \punkte{4}]
Lisa geht zur Schule. Sie läuft \textbf{600\,\einheit{m}} in \textbf{10 Minuten}.

\vspace{0.3em}
\textcolor{orange-warn}{\textbf{Achtung:}} Rechne erst Minuten in Sekunden um!

\vspace{0.5em}

\begin{tabular}{@{}ll}
Umrechnung: & \fz{t} = 10\,\einheit{min} = 10 $\cdot$ 60\,\einheit{s} = \lsg{600}\,\einheit{s} \\[0.8em]
Gegeben: & \fz{s} = \lsg{600\,\einheit{m}}, \quad \fz{t} = \lsg{600\,\einheit{s}} \\[0.5em]
Gesucht: & \lsg{\fz{v}} \\[0.5em]
Rechnung: & \fz{v} = \lsg{\fz{s} / \fz{t} = 600\,\einheit{m} / 600\,\einheit{s}} \\[0.5em]
Ergebnis: & \fz{v} = \lsg{1\,\mps}
\end{tabular}

\vspace{0.3em}
\textcolor{loesung}{\small\textit{Umrechnung 1 BE, Geg+Ges 1 BE, Rechnung 1 BE, Ergebnis 1 BE = 4 BE}}
\end{aufgabenbox}

\vspace{0.8em}

% ═══════════════════════════════════════════════════════════════
% AUFGABE 5: Formel umstellen - Weg (3 BE)
% ═══════════════════════════════════════════════════════════════

\begin{aufgabenbox}[Aufgabe 5: Weg berechnen \sternchen{$\star\star$} \punkte{3}]
Ein Radfahrer fährt mit \fz{v} = 8\,\mps für \fz{t} = 25\,\einheit{s}. \textbf{Berechne} den Weg.

\vspace{0.3em}
\textit{Tipp: Formel umstellen: \fz{s} = \fz{v} $\cdot$ \fz{t}}

\vspace{0.5em}

\begin{tabular}{@{}ll}
Gegeben: & \fz{v} = \lsg{8\,\mps}, \quad \fz{t} = \lsg{25\,\einheit{s}} \\[0.5em]
Gesucht: & \lsg{\fz{s}} \\[0.5em]
Formel: & \fz{s} = \lsg{\fz{v} $\cdot$ \fz{t}} \\[0.5em]
Einsetzen: & \fz{s} = \lsg{8\,\mps $\cdot$ 25\,\einheit{s}} \\[0.5em]
Ergebnis: & \fz{s} = \lsg{200\,\einheit{m}}
\end{tabular}

\vspace{0.3em}
\textcolor{loesung}{\small\textit{Geg+Ges 1 BE, Formel 1 BE, Ergebnis 1 BE = 3 BE}}
\end{aufgabenbox}

\vspace{0.8em}

% ═══════════════════════════════════════════════════════════════
% AUFGABE 6: Formel umstellen - Zeit (3 BE)
% ═══════════════════════════════════════════════════════════════

\begin{aufgabenbox}[Aufgabe 6: Zeit berechnen \sternchen{$\star\star\star$} \punkte{3}]
Ein Bus fährt mit \fz{v} = 20\,\mps eine Strecke von \fz{s} = 400\,\einheit{m}. \textbf{Berechne} die Zeit.

\vspace{0.3em}
\textit{Tipp: Formel umstellen: \fz{t} = \fz{s} / \fz{v}}

\vspace{0.5em}

\begin{tabular}{@{}ll}
Gegeben: & \fz{v} = \lsg{20\,\mps}, \quad \fz{s} = \lsg{400\,\einheit{m}} \\[0.5em]
Gesucht: & \lsg{\fz{t}} \\[0.5em]
Formel: & \fz{t} = \lsg{\fz{s} / \fz{v}} \\[0.5em]
Einsetzen: & \fz{t} = \lsg{400\,\einheit{m} / 20\,\mps} \\[0.5em]
Ergebnis: & \fz{t} = \lsg{20\,\einheit{s}}
\end{tabular}

\vspace{0.3em}
\textcolor{loesung}{\small\textit{Geg+Ges 1 BE, Formel 1 BE, Ergebnis 1 BE = 3 BE}}
\end{aufgabenbox}

\vfill

\begin{center}
\small\textcolor{gray}{Physik Klasse 6 | Geschwindigkeit | ML-02}
\end{center}

\end{document}
